% SGA 3, Expose VI_B, Loop 1 English reconstruction.
% Scope: Theorem 6.2.3 in full, including its proof and source
% editor's note 55.
% Authority: Polo--Gille Exp6B-13oct24.pdf, component-local page 40
% (combined-reader page 381; re-edition header 370).
% Examples 6.2.4 begin on local page 41 / combined page 382 and
% are excluded.
% OCR and Reinhold English Markdown are comparison witnesses only.
% This unit contains no diagrams.

\medskip
\phantomsection\label{thm:sga3-VIB-6.2.3}%
\noindent\textbf{定理 6.2.3.---}
\emph{设 \(S\) 为概形, \(Z\) 为本质自由的 \(S\)-概形,
\(Y\) 为 \(Z\) 的闭子概形.
考虑函子}
\[
  F=\prod_{Z/S}Y/Z:
  (\mathbf{Sch}_{/S})^\circ\longrightarrow\mathbf{Sets}
\]
\emph{其定义如下: 当 \(Y_{S'}\ne Z_{S'}\) 时,
\(F(S')=\varnothing\); 否则 \(F(S')\) 为单元素集.}%
\footnote{\label{note:sga3-VIB-55}编者注: 参见
第~II 篇, \S1, 其中将此函子记为
\(\prod_{Z/S}Y\). 对每个 \(S'\to S\),
\[
  F(S')=\Gamma(Y_{S'}/Z_{S'}),
\]
这里它在 \(Y_{S'}=Z_{S'}\) 时等于
\(\{\operatorname{id}_{Z_{S'}}\}\), 否则为空.
我们还补充了断言 \textup{\SGARefLink{concl:sga3-VIB-6.2.3-i}{(ii)}},
该断言用于 \hyperref[cor:sga3-VIB-6.5.3]{6.5.3} 的证明.}
\begin{enumerate}
\item[\textup{\SGAInlineRefTarget{thm:sga3-VIB-6.2.3:item:i}(i)}]\phantomsection\label{concl:sga3-VIB-6.2.3-i}
  \emph{此函子可由 \(S\) 的一个闭子概形
  \(T\) 表出.}
\item[\textup{(ii)}]\phantomsection\label{concl:sga3-VIB-6.2.3-ii}
  \emph{若此外 \(Y\to Z\) 有限呈示,
  则 \(T\to S\) 也有限呈示.}
\end{enumerate}

\emph{证明.}
首先注意, 所论函子是 fpqc 拓扑下的层.
由于对每个 \(S'\), \(F(S')\) 或为 \(\varnothing\), 或为
\(\{\mathrm{pt}\}\), 这等价于验证
下列断言. 若 \((S_i)\) 为 \(S\) 的开覆盖,
且每个
\[
  Y_{S_i}\longrightarrow Z_{S_i}
\]
都是同构, 则 \(Y\to Z\) 为同构; 而若
\(S'\to S\) 忠实平坦且拟紧, 且
\[
  Y_{S'}\longrightarrow Z_{S'}
\]
为同构, 则 \(Y\to Z\) 仍为同构. 第一个断言
显然成立, 第二个断言由 SGA~1,
VIII, 5.4 或 EGA~IV\(_2\), 2.7.1 得到.

此外, 由 SGA~1, VIII, 1.9, 忠实平坦且
拟紧的态射是闭浸入箭头所成纤维范畴的
有效下降态射.
采用 \hyperref[def:sga3-VIB-6.2.1]{6.2.1} 的记号,
这使我们可以只考虑 \(S=S'_i\) 的情形.

设 \((Z_j)\) 为 \(Z\) 的一个仿射开覆盖,
使得
\[
  \mathcal O(Z_j)
\]
为 \(A=\mathcal O(S)\) 上的投射模. 令
\[
  Y_j=Y\cap Z_j,
\]
并设
\[
  F_j:(\mathbf{Sch}_{/S})^\circ\longrightarrow\mathbf{Sets}
\]
为以 \((Z_j,Y_j)\) 定义的函子, 如同以
\((Z,Y)\) 定义 \(F\) 一样. 它是终函子的子函子,
且显然
\[
  F=\bigcap_j F_j.
\]
因此只需证明每个 \(F_j\) 都可由
\(S\) 的一个闭子概形 \(T_j\) 表出, 因为此时
\(F\) 由 \(S\) 的闭子概形
\[
  T=\bigcap_j T_j
\]
表出.

于是可以假设 \(Z\) 本身仿射, 记
\[
  Z=\operatorname{Spec}B,
\]
其中 \(B\) 为投射 \(A\)-模. 设 \(L\) 为以
\((e_\lambda)_{\lambda\in\Lambda}\) 为基的自由 \(A\)-模,
且 \(B\) 作为 \(A\)-模为其直和因子, 并设
\[
  \phi_\lambda:L\longrightarrow A
\]
为关于此基的坐标线性形式. 设 \(E\)
为 \(B\) 中定义 \(Z\) 的闭子概形 \(Y\) 的理想 \(J\)
的一组生成元, 并设 \(I\) 为 \(A\) 中由坐标
\[
  \phi_\lambda(x),
  \qquad x\in E,\quad \lambda\in\Lambda.
\]
生成的理想.

对每个 \(A\)-代数 \(C\), 态射
\[
  B\otimes_A C\longrightarrow(B/J)\otimes_A C
\]
为同构当且仅当对每个 \(x\in E\), \(x\otimes1\) 在
\(L\otimes_A C\) 中的像为零. 这又
等价于 \(A\to C\) 的核包含
理想 \(I\). 因此
\[
  T=V(I)=\operatorname{Spec}(A/I)
\]
具有所要求的性质, 从而证明了 \textup{\SGARefLink{thm:sga3-VIB-6.2.3:item:i}{(i)}}.

最后, 若 \(Y\to Z\) 有限呈示, 则可选取有限的集合 \(E\).
此时 \(I\) 为 \(A\) 的有限生成理想,
即闭浸入 \(T\to S\) 有限
呈示. 这就证明了 \textup{\SGARefLink{concl:sga3-VIB-6.2.3-i}{(ii)}}.

% 翻译统计: 源行数 130; 输出行数 130 (含本统计行); 行内/行间数学片段 79; 翻译占位 0. 底本疑点另报主线程.

% 独立复核检查点：已逐段对照英文源文，覆盖正文、假设、证明、脚注与源注。
% 行内/陈列数学及全部引用键机械对照通过；仅统一群的invariant为“正规”，保留源文疑点另报。
